\documentclass[11pt,a4paper]{article}
\usepackage[utf8]{inputenc}
\usepackage[T1]{fontenc}
\usepackage[catalan]{babel}
\usepackage{lmodern}
\usepackage[sfdefault,lining]{FiraSans}
\usepackage[a4paper,top=2.3cm,bottom=2.2cm,left=1.8cm,right=1.8cm,headheight=1.3cm,headsep=0.4cm,footskip=1cm]{geometry}
\usepackage{amsmath,amssymb}
\usepackage{icomma}
\usepackage{graphicx}
\usepackage[table]{xcolor}
\usepackage{tikz}
\usepackage{fontawesome5}
\usepackage{ccicons}
\usepackage{fancyhdr}
\usepackage{lastpage}
\usepackage{tasks}
\usepackage{needspace}
\usepackage[most]{tcolorbox}
\usepackage{array}
\usepackage[hidelinks,pdftitle={Fitxa d'exercicis de la unitat 9: Cossos geomètrics. Àrees i volums. 2n d'ESO},pdfsubject={Matemàtiques 2n ESO},pdfkeywords={poliedres, fórmula d'Euler, poliedres regulars, prisma, piràmide, cilindre, con, esfera, àrea, volum, capacitat}]{hyperref}

\definecolor{navy}{HTML}{1B2748}
\definecolor{accent}{HTML}{E8603A}
\definecolor{teal}{HTML}{1F9E90}
\definecolor{soft}{HTML}{EEF1F8}
\definecolor{muted}{HTML}{5D6685}

\setlength{\parindent}{0pt}
\setlength{\parskip}{3pt}
\settasks{label=\alph*), label-format=\bfseries\color{navy}, label-width=1.3em, item-indent=1.8em, column-sep=0.6em, after-item-skip=5pt, before-skip=3pt, after-skip=4pt}

% ---------- capçalera i peu ----------
\newcommand{\llicencia}{\href{https://creativecommons.org/licenses/by-nc-sa/4.0/deed.ca}{\ccbyncsa}\ \ Llicència Creative Commons Reconeixement-NoComercial-CompartirIgual 4.0 Internacional (CC BY-NC-SA 4.0)}
\pagestyle{fancy}
\fancyhf{}
\fancyhead[L]{\raisebox{-0.35\height}{\includegraphics[height=0.95cm]{logo-gva.png}}}
\fancyhead[C]{\small\color{muted}Fitxa d'exercicis · Unitat 9: Cossos geomètrics. Àrees i volums · 2n d'ESO}
\fancyhead[R]{\raisebox{-0.35\height}{\includegraphics[height=0.95cm]{logo-vermellar.png}}}
\fancyfoot[L]{\scriptsize\color{muted}\llicencia}
\fancyfoot[R]{\small\color{muted}\thepage}
\renewcommand{\headrulewidth}{0.4pt}
\renewcommand{\headrule}{\color{muted!40}\hrule height \headrulewidth}
\fancypagestyle{primera}{\fancyhf{}\renewcommand{\headrulewidth}{0pt}
  \fancyfoot[L]{\scriptsize\color{muted}\llicencia}
  \fancyfoot[R]{\small\color{muted}\thepage}}

% ---------- nivells i exercicis ----------
\newcommand{\estrelles}[1]{{\color{accent}\ifcase#1\or\faStar\faStar[regular]\faStar[regular]\or\faStar\faStar\faStar[regular]\or\faStar\faStar\faStar\fi}}
\newcounter{exer}
\newcommand{\ex}[1]{\par\needspace{4\baselineskip}\medskip\stepcounter{exer}\noindent{\color{navy}\bfseries\theexer.}\ \ {\small\estrelles{#1}}\ \ }

\newcommand{\seccio}[2]{\par\needspace{10\baselineskip}\bigskip
  \begin{tcolorbox}[colback=navy,colframe=navy,coltext=white,arc=2mm,boxrule=0pt,left=3mm,right=3mm,top=1.2mm,bottom=1.2mm]
  {\large\bfseries #1\quad #2}
  \end{tcolorbox}\vspace{-2pt}}

\newcolumntype{C}[1]{>{\centering\arraybackslash}p{#1}}
\usepackage{pifont}
% quadre de recordatori (només a la fitxa d'exercicis)
\newcommand{\recorda}[1]{\par\smallskip\begin{tcolorbox}[colback=accent!7,colframe=accent!55,arc=2mm,boxrule=0.5pt,left=2.5mm,right=2.5mm,top=1mm,bottom=1mm]\small{\color{accent}\faLightbulb[regular]}\ \textbf{Recorda:}\ #1\end{tcolorbox}}
\newcommand{\res}[1]{#1}\newcommand{\rest}[1]{#1}

\def\fmneed{14}
% =====================================================================
% Capçalera FentMates de les fitxes: un color per a cada curs.
% Els colors són els de la pàgina inicial (LEVELS a assets/js/materials.js);
% el logo ve de _eines/logo/pdf/ (python3 _eines/logo/genera.py).
%
% Ús, al final del preàmbul de la unitat:
%   % =====================================================================
% Capçalera FentMates de les fitxes: un color per a cada curs.
% Els colors són els de la pàgina inicial (LEVELS a assets/js/materials.js);
% el logo ve de _eines/logo/pdf/ (python3 _eines/logo/genera.py).
%
% Ús, al final del preàmbul de la unitat:
%   % =====================================================================
% Capçalera FentMates de les fitxes: un color per a cada curs.
% Els colors són els de la pàgina inicial (LEVELS a assets/js/materials.js);
% el logo ve de _eines/logo/pdf/ (python3 _eines/logo/genera.py).
%
% Ús, al final del preàmbul de la unitat:
%   \input{../comu/fentmates.tex}
%   \fmcurs{1eso}
% i, al començament del document:
%   \capcalera{Matemàtiques}{1r ESO}{Fitxa d'exercicis · Unitat 2}{Divisibilitat}
% Opcional abans de \input: \def\fmneed{14} (espai mínim abans de cada secció).
% =====================================================================
\makeatletter
\def\fm@logodir{../../logo/pdf}
\providecommand{\fmneed}{10}
\def\fm@cursnom{}

% clau del curs, color i color del text damunt del color
\newcommand{\fm@curs}[3]{\@namedef{fm@color@#1}{#2}\@namedef{fm@text@#1}{#3}}
\fm@curs{1eso}{E6B84A}{navy}
\fm@curs{2eso}{3B82F6}{white}
\fm@curs{3eso}{FF7A50}{white}
\fm@curs{4eso-a}{A78BFA}{white}
\fm@curs{4eso-b}{6D28D9}{white}
\fm@curs{1batx-mi}{9B1C1C}{white}
\fm@curs{1batx-ccss}{EF6B6B}{white}
\fm@curs{2batx-mii}{1E6B47}{white}
\fm@curs{2batx-ccss}{2BB3A3}{white}

\newcommand{\fmcurs}[1]{%
  \@ifundefined{fm@color@#1}{\PackageError{fentmates}{Curs desconegut: #1}{}}{}%
  \edef\fm@tmp{\noexpand\definecolor{curs}{HTML}{\@nameuse{fm@color@#1}}%
               \noexpand\colorlet{curstext}{\@nameuse{fm@text@#1}}}\fm@tmp
  \colorlet{cursfosc}{curs!72!black}%
  \def\fm@logo{\fm@logodir/vertical-#1.pdf}%
  \def\fm@marca{\fm@logodir/marca-#1.pdf}}

% ---------- capçalera de la primera pàgina ----------
% #1 matèria (petit), #2 curs, #3 tipus de document i unitat, #4 títol
\newcommand{\capcalera}[4]{%
  \gdef\fm@cursnom{#2}%
  \par\noindent
  \begin{tcolorbox}[colback=curs,colframe=curs,boxrule=0pt,arc=2.2mm,width=3.5cm,height=2.15cm,
      valign=center,left=3mm,right=2mm,top=1mm,bottom=1mm,nobeforeafter,box align=center]
    \color{curstext}{\footnotesize #1}\par\vspace{1pt}{\fontsize{19}{21}\selectfont\bfseries #2\par}%
  \end{tcolorbox}\hspace{5mm}%
  \begin{minipage}[c]{\dimexpr\linewidth-3.5cm-5mm-2.9cm\relax}
    {\color{cursfosc}\small\bfseries\MakeUppercase{#3}}\par\vspace{2pt}
    {\color{navy}\fontsize{20}{23}\selectfont\bfseries #4\par}
  \end{minipage}\hfill
  \begin{minipage}[c]{2.6cm}\raggedleft\includegraphics[width=2.6cm]{\fm@logo}\end{minipage}\par
  \vspace{2pt}}

% ---------- barres de secció del color del curs ----------
\renewcommand{\seccio}[2]{\par\needspace{\fmneed\baselineskip}\bigskip
  \begin{tcolorbox}[colback=curs,colframe=curs,coltext=curstext,arc=2mm,boxrule=0pt,left=3mm,right=3mm,top=1.2mm,bottom=1.2mm]
  {\large\bfseries #1\quad #2}
  \end{tcolorbox}\vspace{-2pt}}

% ---------- estrelles de dificultat del color del curs (a Batxillerat no se'n posen) ----------
\@ifundefined{estrelles}{}{%
  \renewcommand{\estrelles}[1]{{\color{curs}\ifcase#1\or\faStar\faStar[regular]\faStar[regular]\or\faStar\faStar\faStar[regular]\or\faStar\faStar\faStar\fi}}}

% ---------- capçalera i peu de les pàgines ----------
% primera pàgina: logos institucionals al peu
\fancypagestyle{primera}{\fancyhf{}\renewcommand{\headrulewidth}{0pt}
  \fancyfoot[L]{\raisebox{-0.32\height}{\includegraphics[height=0.78cm]{logo-gva.png}}\hspace{5mm}%
    \raisebox{-0.32\height}{\includegraphics[height=0.78cm]{logo-vermellar.png}}}
  \fancyfoot[R]{\scriptsize\color{muted}\href{https://creativecommons.org/licenses/by-nc-sa/4.0/deed.ca}{\ccbyncsa}\ \ CC BY-NC-SA 4.0 · fentmates.org\qquad\small\thepage}}
% la resta: el símbol a l'esquerra i el curs a la dreta, amb el títol al centre com abans
\fancyhead[L]{\raisebox{-0.3\height}{\includegraphics[height=0.62cm]{\fm@marca}}}
\fancyhead[R]{\small\bfseries\color{cursfosc}\fm@cursnom}
\renewcommand{\headrule}{\color{curs!45}\hrule height \headrulewidth}
\makeatother

%   \fmcurs{1eso}
% i, al començament del document:
%   \capcalera{Matemàtiques}{1r ESO}{Fitxa d'exercicis · Unitat 2}{Divisibilitat}
% Opcional abans de \input: \def\fmneed{14} (espai mínim abans de cada secció).
% =====================================================================
\makeatletter
\def\fm@logodir{../../logo/pdf}
\providecommand{\fmneed}{10}
\def\fm@cursnom{}

% clau del curs, color i color del text damunt del color
\newcommand{\fm@curs}[3]{\@namedef{fm@color@#1}{#2}\@namedef{fm@text@#1}{#3}}
\fm@curs{1eso}{E6B84A}{navy}
\fm@curs{2eso}{3B82F6}{white}
\fm@curs{3eso}{FF7A50}{white}
\fm@curs{4eso-a}{A78BFA}{white}
\fm@curs{4eso-b}{6D28D9}{white}
\fm@curs{1batx-mi}{9B1C1C}{white}
\fm@curs{1batx-ccss}{EF6B6B}{white}
\fm@curs{2batx-mii}{1E6B47}{white}
\fm@curs{2batx-ccss}{2BB3A3}{white}

\newcommand{\fmcurs}[1]{%
  \@ifundefined{fm@color@#1}{\PackageError{fentmates}{Curs desconegut: #1}{}}{}%
  \edef\fm@tmp{\noexpand\definecolor{curs}{HTML}{\@nameuse{fm@color@#1}}%
               \noexpand\colorlet{curstext}{\@nameuse{fm@text@#1}}}\fm@tmp
  \colorlet{cursfosc}{curs!72!black}%
  \def\fm@logo{\fm@logodir/vertical-#1.pdf}%
  \def\fm@marca{\fm@logodir/marca-#1.pdf}}

% ---------- capçalera de la primera pàgina ----------
% #1 matèria (petit), #2 curs, #3 tipus de document i unitat, #4 títol
\newcommand{\capcalera}[4]{%
  \gdef\fm@cursnom{#2}%
  \par\noindent
  \begin{tcolorbox}[colback=curs,colframe=curs,boxrule=0pt,arc=2.2mm,width=3.5cm,height=2.15cm,
      valign=center,left=3mm,right=2mm,top=1mm,bottom=1mm,nobeforeafter,box align=center]
    \color{curstext}{\footnotesize #1}\par\vspace{1pt}{\fontsize{19}{21}\selectfont\bfseries #2\par}%
  \end{tcolorbox}\hspace{5mm}%
  \begin{minipage}[c]{\dimexpr\linewidth-3.5cm-5mm-2.9cm\relax}
    {\color{cursfosc}\small\bfseries\MakeUppercase{#3}}\par\vspace{2pt}
    {\color{navy}\fontsize{20}{23}\selectfont\bfseries #4\par}
  \end{minipage}\hfill
  \begin{minipage}[c]{2.6cm}\raggedleft\includegraphics[width=2.6cm]{\fm@logo}\end{minipage}\par
  \vspace{2pt}}

% ---------- barres de secció del color del curs ----------
\renewcommand{\seccio}[2]{\par\needspace{\fmneed\baselineskip}\bigskip
  \begin{tcolorbox}[colback=curs,colframe=curs,coltext=curstext,arc=2mm,boxrule=0pt,left=3mm,right=3mm,top=1.2mm,bottom=1.2mm]
  {\large\bfseries #1\quad #2}
  \end{tcolorbox}\vspace{-2pt}}

% ---------- estrelles de dificultat del color del curs (a Batxillerat no se'n posen) ----------
\@ifundefined{estrelles}{}{%
  \renewcommand{\estrelles}[1]{{\color{curs}\ifcase#1\or\faStar\faStar[regular]\faStar[regular]\or\faStar\faStar\faStar[regular]\or\faStar\faStar\faStar\fi}}}

% ---------- capçalera i peu de les pàgines ----------
% primera pàgina: logos institucionals al peu
\fancypagestyle{primera}{\fancyhf{}\renewcommand{\headrulewidth}{0pt}
  \fancyfoot[L]{\raisebox{-0.32\height}{\includegraphics[height=0.78cm]{logo-gva.png}}\hspace{5mm}%
    \raisebox{-0.32\height}{\includegraphics[height=0.78cm]{logo-vermellar.png}}}
  \fancyfoot[R]{\scriptsize\color{muted}\href{https://creativecommons.org/licenses/by-nc-sa/4.0/deed.ca}{\ccbyncsa}\ \ CC BY-NC-SA 4.0 · fentmates.org\qquad\small\thepage}}
% la resta: el símbol a l'esquerra i el curs a la dreta, amb el títol al centre com abans
\fancyhead[L]{\raisebox{-0.3\height}{\includegraphics[height=0.62cm]{\fm@marca}}}
\fancyhead[R]{\small\bfseries\color{cursfosc}\fm@cursnom}
\renewcommand{\headrule}{\color{curs!45}\hrule height \headrulewidth}
\makeatother

%   \fmcurs{1eso}
% i, al començament del document:
%   \capcalera{Matemàtiques}{1r ESO}{Fitxa d'exercicis · Unitat 2}{Divisibilitat}
% Opcional abans de \input: \def\fmneed{14} (espai mínim abans de cada secció).
% =====================================================================
\makeatletter
\def\fm@logodir{../../logo/pdf}
\providecommand{\fmneed}{10}
\def\fm@cursnom{}

% clau del curs, color i color del text damunt del color
\newcommand{\fm@curs}[3]{\@namedef{fm@color@#1}{#2}\@namedef{fm@text@#1}{#3}}
\fm@curs{1eso}{E6B84A}{navy}
\fm@curs{2eso}{3B82F6}{white}
\fm@curs{3eso}{FF7A50}{white}
\fm@curs{4eso-a}{A78BFA}{white}
\fm@curs{4eso-b}{6D28D9}{white}
\fm@curs{1batx-mi}{9B1C1C}{white}
\fm@curs{1batx-ccss}{EF6B6B}{white}
\fm@curs{2batx-mii}{1E6B47}{white}
\fm@curs{2batx-ccss}{2BB3A3}{white}

\newcommand{\fmcurs}[1]{%
  \@ifundefined{fm@color@#1}{\PackageError{fentmates}{Curs desconegut: #1}{}}{}%
  \edef\fm@tmp{\noexpand\definecolor{curs}{HTML}{\@nameuse{fm@color@#1}}%
               \noexpand\colorlet{curstext}{\@nameuse{fm@text@#1}}}\fm@tmp
  \colorlet{cursfosc}{curs!72!black}%
  \def\fm@logo{\fm@logodir/vertical-#1.pdf}%
  \def\fm@marca{\fm@logodir/marca-#1.pdf}}

% ---------- capçalera de la primera pàgina ----------
% #1 matèria (petit), #2 curs, #3 tipus de document i unitat, #4 títol
\newcommand{\capcalera}[4]{%
  \gdef\fm@cursnom{#2}%
  \par\noindent
  \begin{tcolorbox}[colback=curs,colframe=curs,boxrule=0pt,arc=2.2mm,width=3.5cm,height=2.15cm,
      valign=center,left=3mm,right=2mm,top=1mm,bottom=1mm,nobeforeafter,box align=center]
    \color{curstext}{\footnotesize #1}\par\vspace{1pt}{\fontsize{19}{21}\selectfont\bfseries #2\par}%
  \end{tcolorbox}\hspace{5mm}%
  \begin{minipage}[c]{\dimexpr\linewidth-3.5cm-5mm-2.9cm\relax}
    {\color{cursfosc}\small\bfseries\MakeUppercase{#3}}\par\vspace{2pt}
    {\color{navy}\fontsize{20}{23}\selectfont\bfseries #4\par}
  \end{minipage}\hfill
  \begin{minipage}[c]{2.6cm}\raggedleft\includegraphics[width=2.6cm]{\fm@logo}\end{minipage}\par
  \vspace{2pt}}

% ---------- barres de secció del color del curs ----------
\renewcommand{\seccio}[2]{\par\needspace{\fmneed\baselineskip}\bigskip
  \begin{tcolorbox}[colback=curs,colframe=curs,coltext=curstext,arc=2mm,boxrule=0pt,left=3mm,right=3mm,top=1.2mm,bottom=1.2mm]
  {\large\bfseries #1\quad #2}
  \end{tcolorbox}\vspace{-2pt}}

% ---------- estrelles de dificultat del color del curs (a Batxillerat no se'n posen) ----------
\@ifundefined{estrelles}{}{%
  \renewcommand{\estrelles}[1]{{\color{curs}\ifcase#1\or\faStar\faStar[regular]\faStar[regular]\or\faStar\faStar\faStar[regular]\or\faStar\faStar\faStar\fi}}}

% ---------- capçalera i peu de les pàgines ----------
% primera pàgina: logos institucionals al peu
\fancypagestyle{primera}{\fancyhf{}\renewcommand{\headrulewidth}{0pt}
  \fancyfoot[L]{\raisebox{-0.32\height}{\includegraphics[height=0.78cm]{logo-gva.png}}\hspace{5mm}%
    \raisebox{-0.32\height}{\includegraphics[height=0.78cm]{logo-vermellar.png}}}
  \fancyfoot[R]{\scriptsize\color{muted}\href{https://creativecommons.org/licenses/by-nc-sa/4.0/deed.ca}{\ccbyncsa}\ \ CC BY-NC-SA 4.0 · fentmates.org\qquad\small\thepage}}
% la resta: el símbol a l'esquerra i el curs a la dreta, amb el títol al centre com abans
\fancyhead[L]{\raisebox{-0.3\height}{\includegraphics[height=0.62cm]{\fm@marca}}}
\fancyhead[R]{\small\bfseries\color{cursfosc}\fm@cursnom}
\renewcommand{\headrule}{\color{curs!45}\hrule height \headrulewidth}
\makeatother

\fmcurs{2eso}

\begin{document}
\thispagestyle{primera}

\capcalera{Matemàtiques}{2n ESO}{Fitxa d'exercicis · Unitat 9}{Cossos geomètrics.\\ Àrees i volums}
\vspace{10pt}
Nom i cognoms: \rule{8.4cm}{0.4pt}\hfill Grup: \rule{1.6cm}{0.4pt}\hfill Data: \rule{2.2cm}{0.4pt}
\vspace{8pt}

\begin{tcolorbox}[colback=soft,colframe=soft,arc=2mm,left=3mm,right=3mm,top=1.5mm,bottom=1.5mm]
\small Exercicis graduats per nivells de dificultat:\quad \estrelles{1}\ bàsic \qquad \estrelles{2}\ mitjà \qquad \estrelles{3}\ avançat.\\[2pt] Resol-los al quadern mostrant tots els passos i fes un dibuix de cada cos. Si un resultat no és exacte, arrodoneix-lo a les centèsimes i fes servir la tecla $\pi$ de la calculadora.
\end{tcolorbox}

% =====================================================
\seccio{1}{Poliedres. Fórmula d'Euler. Poliedres regulars}

\ex{1} Copia i completa la taula. Comprova en cada cas la fórmula d'Euler.\par\smallskip\begin{center}\renewcommand{\arraystretch}{1.25}\begin{tabular}{|l|C{1.6cm}|C{1.6cm}|C{1.6cm}|}\hline Poliedre & Cares & Vèrtexs & Arestes \\\hline prisma pentagonal & & &  \\\hline piràmide hexagonal & & &  \\\hline prisma triangular & & &  \\\hline piràmide quadrangular & & &  \\\hline ortoedre & & &  \\\hline\end{tabular}\end{center}\recorda{En un poliedre convex, \textbf{cares} + \textbf{vèrtexs} = \textbf{arestes} + 2: \ $C+V=A+2$ (fórmula d'Euler). Un prisma de base de $n$ costats té $n+2$ cares, $2n$ vèrtexs i $3n$ arestes; una piràmide, $n+1$ cares, $n+1$ vèrtexs i $2n$ arestes.}

\ex{1} Fes servir la fórmula d'Euler per a calcular el que falta en cada poliedre convex:\begin{tasks}(3)
  \task $C=6$, \ $V=8$, \ $A=\ ?$
  \task $V=12$, \ $A=18$, \ $C=\ ?$
  \task $C=7$, \ $A=12$, \ $V=\ ?$
\end{tasks}

\ex{1} Quins d'aquests poliedres són regulars? Justifica-ho.\begin{tasks}(3)
  \task tetraedre
  \task prisma triangular
  \task cub
  \task piràmide quadrangular
  \task dodecaedre
  \task ortoedre
\end{tasks}

\ex{2} Escriu, per a cadascun dels cinc poliedres regulars, quantes cares té, quina forma tenen i quantes cares es troben a cada vèrtex.

\ex{2} Un prisma té 18 arestes. Quin polígon té per base? Quantes cares i quants vèrtexs té?

\ex{2} Una piràmide té 9 cares. Quin polígon té per base? Quants vèrtexs i arestes té?

\ex{3} Pot existir un poliedre convex amb: \ a) 8 cares, 12 vèrtexs i 18 arestes? \ b) 6 cares, 7 vèrtexs i 12 arestes?

\ex{3} Explica per què no hi pot haver cap poliedre regular amb cares hexagonals.


% =====================================================
\seccio{2}{Prismes: elements i àrea}

\ex{1} \begin{minipage}[t]{0.50\linewidth}Calcula l'àrea total:\begin{tasks}(1)
  \task d'un cub de 5 cm d'aresta;
  \task de l'ortoedre de la figura.
\end{tasks}\end{minipage}\hfill\begin{minipage}[t]{0.36\linewidth}\vspace{0pt}\centering \begin{tikzpicture}[x=0.42cm,y=0.42cm,baseline=(current bounding box.center),line join=round]
\path[fill=teal!14,draw=none] (0.000,2.782) -- (0.000,0.000) -- (6.181,-1.231) -- (6.181,1.550) -- cycle;
\path[fill=teal!14,draw=none] (6.181,1.550) -- (6.181,-1.231) -- (8.059,0.092) -- (8.059,2.874) -- cycle;
\path[fill=teal!14,draw=none] (1.878,4.105) -- (0.000,2.782) -- (6.181,1.550) -- (8.059,2.874) -- cycle;
\draw[dashed,draw=muted] (0.000,0.000) -- (1.878,1.323);
\draw[dashed,draw=muted] (1.878,1.323) -- (8.059,0.092);
\draw[dashed,draw=muted] (1.878,1.323) -- (1.878,4.105);
\draw[thick] (0.000,0.000) -- (6.181,-1.231);
\draw[thick] (6.181,-1.231) -- (8.059,0.092);
\draw[thick] (0.000,0.000) -- (0.000,2.782);
\draw[thick] (6.181,-1.231) -- (6.181,1.550);
\draw[thick] (0.000,2.782) -- (6.181,1.550);
\draw[thick] (0.000,2.782) -- (1.878,4.105);
\draw[thick] (8.059,0.092) -- (8.059,2.874);
\draw[thick] (6.181,1.550) -- (8.059,2.874);
\draw[thick] (1.878,4.105) -- (8.059,2.874);
\node[font=\small,inner sep=2pt,below] at (3.090,-0.616) {$7$\,cm};
\node[font=\small,inner sep=2pt,right] at (7.120,-0.570) {$4$\,cm};
\node[font=\small,inner sep=2pt,right] at (6.181,0.160) {$3$\,cm};
\end{tikzpicture}\end{minipage}\par\recorda{Àrea d'un prisma recte: $A_L=P_{\text{base}}\cdot h$, \ $A_T=A_L+2\,A_{\text{base}}$. Ortoedre d'arestes $a,\,b,\,c$: $A_T=2(ab+ac+bc)$ i la diagonal és $D=\sqrt{a^2+b^2+c^2}$.}

\ex{1} Calcula la diagonal d'un ortoedre de 6 cm × 2 cm × 3 cm.

\ex{1} Un prisma quadrangular regular té 4 cm d'aresta de la base i 10 cm d'altura. Calcula'n l'àrea lateral i l'àrea total.

\ex{2} La base d'un prisma recte és un triangle rectangle de catets 5 cm i 12 cm, i l'altura del prisma és de 10 cm. Calcula'n l'àrea total.

\ex{2} Calcula l'àrea total d'un prisma hexagonal regular de 8 cm d'aresta de la base i 5 cm d'altura.

\ex{2} Volem fer una caixa sense tapa de 30 cm × 20 cm de base i 15 cm d'altura. Quant cartó necessitem?

\ex{3} Un cub té una àrea total de 216 cm². Quant mesura l'aresta? I la diagonal?

\ex{3} Volem pintar les quatre parets d'una habitació de 4 m × 3,5 m i 2,5 m d'altura. La porta fa 0,8 m × 2 m i la finestra 1,2 m × 1 m. Si pintar costa 4 € el metre quadrat, quant ens costarà?


% =====================================================
\seccio{3}{Piràmides: elements i àrea}

\ex{1} \begin{minipage}[t]{0.54\linewidth}La piràmide quadrangular regular de la figura té 16 cm d'aresta de la base i 6 cm d'altura. Calcula'n l'apotema i l'àrea total.\end{minipage}\hfill\begin{minipage}[t]{0.32\linewidth}\vspace{0pt}\centering \begin{tikzpicture}[x=0.6cm,y=0.6cm,baseline=(current bounding box.center),line join=round]
\path[fill=teal!14,draw=none] (0.000,2.967) -- (-2.434,-0.279) -- (0.744,-0.912) -- cycle;
\path[fill=teal!14,draw=none] (0.000,2.967) -- (0.744,-0.912) -- (2.434,0.279) -- cycle;
\draw[dashed,draw=muted] (-2.434,-0.279) -- (-0.744,0.912);
\draw[dashed,draw=muted] (2.434,0.279) -- (-0.744,0.912);
\draw[dashed,draw=muted] (-0.744,0.912) -- (0.000,2.967);
\draw[thick] (-2.434,-0.279) -- (0.744,-0.912);
\draw[thick] (0.744,-0.912) -- (2.434,0.279);
\draw[thick] (-2.434,-0.279) -- (0.000,2.967);
\draw[thick] (0.744,-0.912) -- (0.000,2.967);
\draw[thick] (2.434,0.279) -- (0.000,2.967);
\node[font=\small,inner sep=2pt,below] at (-0.845,-0.595) {$16$\,cm};
\draw[dashed,draw=accent,thick] (0.000,0.000) -- (0.000,2.967);
\draw[dashed,draw=accent,thick] (0.000,0.000) -- (1.589,-0.317);
\draw[draw=navy,thick] (1.589,-0.317) -- (0.000,2.967);
\fill (0.000,0.000) circle (1.3pt);
\node[font=\small,inner sep=2pt,left,text=accent] at (0.000,1.483) {$6$\,cm};
\end{tikzpicture}\end{minipage}\par\recorda{En una piràmide regular, l'altura $h$, l'apotema de la base $a_b$ i l'apotema de la piràmide $a_p$ formen un triangle rectangle: $a_p^2=h^2+a_b^2$. \ $A_L=\frac{P_{\text{base}}\cdot a_p}{2}$, \ $A_T=A_L+A_{\text{base}}$.}

\ex{1} Una piràmide hexagonal regular té 6 cm d'aresta de la base i 10 cm d'apotema. Calcula'n l'àrea lateral.

\ex{1} Una piràmide triangular regular té 8 cm d'aresta de la base i 12 cm d'apotema. Calcula'n l'àrea lateral.

\ex{2} Una piràmide quadrangular regular té 10 cm d'aresta de la base i 13 cm d'apotema. Calcula'n l'altura i l'àrea total.

\ex{2} Una piràmide quadrangular regular té 6 cm d'aresta de la base i 4 cm d'altura. Calcula l'apotema, l'aresta lateral i l'àrea total.

\ex{2} Calcula l'àrea total d'un tetraedre regular de 4 cm d'aresta.

\ex{3} Una tenda de campanya té forma de piràmide quadrangular regular de 2 m de costat de la base i 1,5 m d'altura. Quanta lona cal per a fer-la, si no té terra?

\ex{3} Calcula l'àrea total d'una piràmide hexagonal regular de 4 cm d'aresta de la base i 6 cm d'altura.


% =====================================================
\seccio{4}{Cossos de revolució: cilindre i con}

\ex{1} Quin cos de revolució s'obté en fer girar:\begin{tasks}(1)
  \task un rectangle al voltant d'un costat?
  \task un triangle rectangle al voltant d'un catet?
  \task un semicercle al voltant del diàmetre?
  \task un triangle rectangle al voltant de la hipotenusa?
\end{tasks}

\ex{1} \begin{minipage}[t]{0.60\linewidth}Calcula l'àrea lateral i l'àrea total del cilindre de la figura, de 3 cm de radi i 8 cm d'altura.\end{minipage}\hfill\begin{minipage}[t]{0.26\linewidth}\vspace{0pt}\centering \begin{tikzpicture}[x=0.45cm,y=0.45cm,baseline=(current bounding box.center),line join=round]
\path[fill=teal!14,draw=none] (-1.500,3.709) -- (-1.500,0.000) -- (-1.497,-0.037) -- (-1.487,-0.073) -- (-1.471,-0.110) -- (-1.449,-0.145) -- (-1.420,-0.181) -- (-1.386,-0.215) -- (-1.345,-0.249) -- (-1.299,-0.281) -- (-1.247,-0.312) -- (-1.190,-0.342) -- (-1.128,-0.370) -- (-1.061,-0.397) -- (-0.989,-0.422) -- (-0.913,-0.446) -- (-0.833,-0.467) -- (-0.750,-0.487) -- (-0.663,-0.504) -- (-0.574,-0.519) -- (-0.482,-0.532) -- (-0.388,-0.543) -- (-0.293,-0.551) -- (-0.196,-0.557) -- (-0.098,-0.561) -- (-0.000,-0.562) -- (0.098,-0.561) -- (0.196,-0.557) -- (0.293,-0.551) -- (0.388,-0.543) -- (0.482,-0.532) -- (0.574,-0.519) -- (0.663,-0.504) -- (0.750,-0.487) -- (0.833,-0.467) -- (0.913,-0.446) -- (0.989,-0.422) -- (1.061,-0.397) -- (1.128,-0.370) -- (1.190,-0.342) -- (1.247,-0.312) -- (1.299,-0.281) -- (1.345,-0.249) -- (1.386,-0.215) -- (1.420,-0.181) -- (1.449,-0.145) -- (1.471,-0.110) -- (1.487,-0.073) -- (1.497,-0.037) -- (1.500,-0.000) -- (1.500,3.709) -- (1.500,3.709) -- (1.497,3.672) -- (1.487,3.635) -- (1.471,3.599) -- (1.449,3.563) -- (1.420,3.528) -- (1.386,3.494) -- (1.345,3.460) -- (1.299,3.428) -- (1.247,3.397) -- (1.190,3.367) -- (1.128,3.338) -- (1.061,3.311) -- (0.989,3.286) -- (0.913,3.263) -- (0.833,3.242) -- (0.750,3.222) -- (0.663,3.205) -- (0.574,3.190) -- (0.482,3.177) -- (0.388,3.166) -- (0.293,3.158) -- (0.196,3.152) -- (0.098,3.148) -- (-0.000,3.147) -- (-0.098,3.148) -- (-0.196,3.152) -- (-0.293,3.158) -- (-0.388,3.166) -- (-0.482,3.177) -- (-0.574,3.190) -- (-0.663,3.205) -- (-0.750,3.222) -- (-0.833,3.242) -- (-0.913,3.263) -- (-0.989,3.286) -- (-1.061,3.311) -- (-1.128,3.338) -- (-1.190,3.367) -- (-1.247,3.397) -- (-1.299,3.428) -- (-1.345,3.460) -- (-1.386,3.494) -- (-1.420,3.528) -- (-1.449,3.563) -- (-1.471,3.599) -- (-1.487,3.635) -- (-1.497,3.672) -- (-1.500,3.709) -- cycle;
\draw[dashed,draw=muted] (1.500,0.000) -- (1.497,0.037) -- (1.487,0.073) -- (1.471,0.110) -- (1.449,0.145) -- (1.420,0.181) -- (1.386,0.215) -- (1.345,0.249) -- (1.299,0.281) -- (1.247,0.312) -- (1.190,0.342) -- (1.128,0.370) -- (1.061,0.397) -- (0.989,0.422) -- (0.913,0.446) -- (0.833,0.467) -- (0.750,0.487) -- (0.663,0.504) -- (0.574,0.519) -- (0.482,0.532) -- (0.388,0.543) -- (0.293,0.551) -- (0.196,0.557) -- (0.098,0.561) -- (0.000,0.562) -- (-0.098,0.561) -- (-0.196,0.557) -- (-0.293,0.551) -- (-0.388,0.543) -- (-0.482,0.532) -- (-0.574,0.519) -- (-0.663,0.504) -- (-0.750,0.487) -- (-0.833,0.467) -- (-0.913,0.446) -- (-0.989,0.422) -- (-1.061,0.397) -- (-1.128,0.370) -- (-1.190,0.342) -- (-1.247,0.312) -- (-1.299,0.281) -- (-1.345,0.249) -- (-1.386,0.215) -- (-1.420,0.181) -- (-1.449,0.145) -- (-1.471,0.110) -- (-1.487,0.073) -- (-1.497,0.037) -- (-1.500,0.000);
\draw[thick] (-1.500,0.000) -- (-1.497,-0.037) -- (-1.487,-0.073) -- (-1.471,-0.110) -- (-1.449,-0.145) -- (-1.420,-0.181) -- (-1.386,-0.215) -- (-1.345,-0.249) -- (-1.299,-0.281) -- (-1.247,-0.312) -- (-1.190,-0.342) -- (-1.128,-0.370) -- (-1.061,-0.397) -- (-0.989,-0.422) -- (-0.913,-0.446) -- (-0.833,-0.467) -- (-0.750,-0.487) -- (-0.663,-0.504) -- (-0.574,-0.519) -- (-0.482,-0.532) -- (-0.388,-0.543) -- (-0.293,-0.551) -- (-0.196,-0.557) -- (-0.098,-0.561) -- (-0.000,-0.562) -- (0.098,-0.561) -- (0.196,-0.557) -- (0.293,-0.551) -- (0.388,-0.543) -- (0.482,-0.532) -- (0.574,-0.519) -- (0.663,-0.504) -- (0.750,-0.487) -- (0.833,-0.467) -- (0.913,-0.446) -- (0.989,-0.422) -- (1.061,-0.397) -- (1.128,-0.370) -- (1.190,-0.342) -- (1.247,-0.312) -- (1.299,-0.281) -- (1.345,-0.249) -- (1.386,-0.215) -- (1.420,-0.181) -- (1.449,-0.145) -- (1.471,-0.110) -- (1.487,-0.073) -- (1.497,-0.037) -- (1.500,-0.000);
\draw[thick] (-1.500,0.000) -- (-1.500,3.709);
\draw[thick] (1.500,0.000) -- (1.500,3.709);
\draw[fill=teal!26,draw=black,thick] (1.500,3.709) -- (1.487,3.782) -- (1.449,3.854) -- (1.386,3.924) -- (1.299,3.990) -- (1.190,4.051) -- (1.061,4.106) -- (0.913,4.155) -- (0.750,4.195) -- (0.574,4.228) -- (0.388,4.251) -- (0.196,4.266) -- (0.000,4.271) -- (-0.196,4.266) -- (-0.388,4.251) -- (-0.574,4.228) -- (-0.750,4.195) -- (-0.913,4.155) -- (-1.061,4.106) -- (-1.190,4.051) -- (-1.299,3.990) -- (-1.386,3.924) -- (-1.449,3.854) -- (-1.487,3.782) -- (-1.500,3.709) -- (-1.487,3.635) -- (-1.449,3.563) -- (-1.386,3.494) -- (-1.299,3.428) -- (-1.190,3.367) -- (-1.061,3.311) -- (-0.913,3.263) -- (-0.750,3.222) -- (-0.574,3.190) -- (-0.388,3.166) -- (-0.196,3.152) -- (-0.000,3.147) -- (0.196,3.152) -- (0.388,3.166) -- (0.574,3.190) -- (0.750,3.222) -- (0.913,3.263) -- (1.061,3.311) -- (1.190,3.367) -- (1.299,3.428) -- (1.386,3.494) -- (1.449,3.563) -- (1.487,3.635) -- (1.500,3.709) -- cycle;
\draw[dashed,draw=accent,thick] (0.000,0.000) -- (1.500,0.000);
\fill (0.000,0.000) circle (1.3pt);
\node[font=\small,inner sep=2pt,above,text=accent] at (0.750,0.000) {$3$\,cm};
\node[font=\small,inner sep=2pt,right] at (1.500,1.854) {$8$\,cm};
\end{tikzpicture}\end{minipage}\par\recorda{Cilindre: $A_L=2\pi r h$, \ $A_T=2\pi r h+2\pi r^2$. \ Con: $A_L=\pi r g$, \ $A_T=\pi r g+\pi r^2$, on la generatriu compleix $g^2=h^2+r^2$.}

\ex{1} Un con té 6 cm de radi i 10 cm de generatriu. Calcula'n l'àrea lateral i l'àrea total.

\ex{2} Un con té 16 cm de diàmetre i 15 cm d'altura. Calcula la generatriu i l'àrea total.

\ex{2} L'etiqueta d'una llanda cilíndrica de 8 cm de diàmetre i 11 cm d'altura cobreix tota la superfície lateral. Quina àrea té?

\ex{2} Un con té 25 cm de generatriu i 7 cm de radi. Quina altura té?

\ex{3} Un cilindre de 10 cm d'altura té una àrea lateral de 314,16 cm². Quin radi té? Calcula'n l'àrea total.

\ex{3} Volem fer gots cilíndrics de cartolina, sense tapa, de 3 cm de radi i 9 cm d'altura. Quanta cartolina cal per a cada got? Quants gots, com a màxim, podríem fer amb 1 m² de cartolina (sense comptar el que es perd en retallar)?


% =====================================================
\seccio{5}{L'esfera i la seua àrea}

\ex{1} Calcula l'àrea de:\begin{tasks}(2)
  \task una esfera de 7 cm de radi;
  \task una esfera de 10 cm de diàmetre.
\end{tasks}\recorda{L'àrea de l'esfera és la de quatre cercles màxims: $A=4\pi r^2$. Compte: si et donen el diàmetre, el radi n'és la meitat.}

\ex{1} Una esfera té una àrea de $36\pi$ cm². Quin radi té?

\ex{2} Una pilota de tennis fa 6,5 cm de diàmetre. Quanta superfície té?

\ex{2} Volem pintar per fora una cúpula semiesfèrica de 6 m de radi. Quina superfície hem de pintar?

\ex{2} La circumferència màxima d'una pilota fa 50 cm. Calcula'n el radi i l'àrea.

\ex{3} Una esfera està ficada dins d'un cub de 10 cm d'aresta i toca totes les cares. Quina és l'àrea de l'esfera? Quin percentatge representa de l'àrea del cub?

\ex{3} Una esfera de 4 cm de radi està ficada en un cilindre que l'envolta exactament (radi 4 cm i altura 8 cm). Comprova que l'àrea de l'esfera és igual a l'àrea lateral del cilindre. Quina fracció de l'àrea total del cilindre representa?


% =====================================================
\seccio{6}{Volum. Prismes i cilindres}

\ex{1} Expressa en la unitat que s'indica:\begin{tasks}(3)
  \task 5 m³ \ en dm³
  \task 2,5 dm³ \ en cm³
  \task 7500 cm³ \ en dm³
  \task 0,35 m³ \ en litres
  \task 3 L \ en cm³
  \task 250 mL \ en dm³
\end{tasks}\recorda{Cada unitat de volum és 1000 vegades la següent: 1 m³ = 1000 dm³, 1 dm³ = 1000 cm³. \ I la clau de la capacitat: 1 dm³ = 1 L, \ 1 cm³ = 1 mL, \ 1 m³ = 1000 L.}

\ex{1} Calcula el volum d'un cub de 6 cm d'aresta i el d'un ortoedre de 10 cm × 4 cm × 5 cm.

\ex{1} Calcula el volum d'un cilindre de 2 cm de radi i 10 cm d'altura.

\ex{2} Un prisma té de base un triangle rectangle de catets 3 cm i 4 cm, i 20 cm d'altura. Quin volum té?

\ex{2} Un aquari fa 60 cm de llarg, 30 cm d'ample i 40 cm d'alt. Quants litres hi caben? Quants n'hi posem si l'omplim fins a 35 cm d'altura?

\ex{2} Un dipòsit cilíndric té 1 m de radi i 3 m d'altura. Quants litres d'aigua hi caben?

\ex{3} Calcula el volum d'un prisma hexagonal regular de 2 cm d'aresta de la base i 10 cm d'altura.

\ex{3} Un tetrabric d'1 L de llet té una base rectangular de 6 cm × 9,5 cm. Quina altura ha de tenir com a mínim?


% =====================================================
\seccio{7}{Volum de piràmides, cons i esferes}

\ex{1} Calcula el volum d'una piràmide quadrangular de 9 cm d'aresta de la base i 10 cm d'altura.\recorda{$V_{\text{piràmide}}=\frac{A_{\text{base}}\cdot h}{3}$, \ $V_{\text{con}}=\frac{\pi r^2h}{3}$ (un terç del prisma o del cilindre de la mateixa base i altura), \ $V_{\text{esfera}}=\frac43\pi r^3$.}

\ex{1} Calcula el volum d'un con de 4 cm de radi i 9 cm d'altura.

\ex{1} Calcula el volum d'una esfera de 3 cm de radi i el d'una de 10 cm de radi.

\ex{2} Un con té 15 cm de generatriu i 9 cm de radi. Calcula'n l'altura i el volum.

\ex{2} Una piràmide quadrangular regular té 8 cm d'aresta de la base i 5 cm d'apotema. Calcula'n l'altura i el volum.

\ex{2} Una pilota de bàsquet fa 24 cm de diàmetre. Quants litres d'aire conté?

\ex{3} Una esfera té un volum de $288\pi$ cm³. Quin radi té? Calcula'n l'àrea.

\ex{3} Un cilindre de 3 cm de radi i 6 cm d'altura està ple d'aigua. Hi fiquem una esfera de 3 cm de radi (que hi cap justa). Quina fracció de l'aigua vessa?


% =====================================================
\seccio{8}{Cossos compostos i problemes}

\ex{1} \begin{minipage}[t]{0.66\linewidth}La peça de la figura és un cub de 5 cm d'aresta amb una piràmide de 6 cm d'altura al damunt. Calcula'n el volum.\end{minipage}\hfill\begin{minipage}[t]{0.2\linewidth}\vspace{0pt}\centering \begin{tikzpicture}[x=0.3cm,y=0.3cm,baseline=(current bounding box.center),line join=round]
\path[fill=teal!14,draw=none] (0.000,4.636) -- (0.000,0.000) -- (4.415,-0.879) -- (4.415,3.757) -- cycle;
\path[fill=teal!14,draw=none] (4.415,3.757) -- (4.415,-0.879) -- (6.762,0.774) -- (6.762,5.410) -- cycle;
\path[fill=teal!14,draw=none] (3.381,10.586) -- (0.000,4.636) -- (4.415,3.757) -- cycle;
\path[fill=teal!14,draw=none] (3.381,10.586) -- (4.415,3.757) -- (6.762,5.410) -- cycle;
\draw[dashed,draw=muted] (0.000,0.000) -- (2.347,1.654);
\draw[dashed,draw=muted] (2.347,1.654) -- (6.762,0.774);
\draw[dashed,draw=muted] (2.347,1.654) -- (2.347,6.290);
\draw[dashed,draw=muted] (0.000,4.636) -- (2.347,6.290);
\draw[dashed,draw=muted] (2.347,6.290) -- (6.762,5.410);
\draw[dashed,draw=muted] (2.347,6.290) -- (3.381,10.586);
\draw[thick] (0.000,0.000) -- (4.415,-0.879);
\draw[thick] (4.415,-0.879) -- (6.762,0.774);
\draw[thick] (0.000,0.000) -- (0.000,4.636);
\draw[thick] (4.415,-0.879) -- (4.415,3.757);
\draw[thick] (0.000,4.636) -- (4.415,3.757);
\draw[thick] (6.762,0.774) -- (6.762,5.410);
\draw[thick] (4.415,3.757) -- (6.762,5.410);
\draw[thick] (0.000,4.636) -- (3.381,10.586);
\draw[thick] (4.415,3.757) -- (3.381,10.586);
\draw[thick] (6.762,5.410) -- (3.381,10.586);
\end{tikzpicture}\end{minipage}\par\recorda{Per a calcular el volum d'un cos compost, el descomponem en cossos simples i sumem els volums (o els restem, si hi ha forats). Per a l'àrea, només sumem les superfícies que queden per fora.}

\ex{2} Un gelat és un con de 3 cm de radi i 12 cm d'altura, ple de gelat, amb una semiesfera de gelat al damunt. Quant de gelat hi ha?

\ex{2} Una sitja és un cilindre de 3 m de radi i 8 m d'altura amb una teulada cònica de 4 m d'altura. Quin volum té?

\ex{2} Un tub de 2 m de llarg té 6 cm de radi exterior i 5 cm de radi interior. Quin volum de material té?

\ex{2} Un cub de 12 cm d'aresta té un forat cilíndric de 2 cm de radi que el travessa de banda a banda. Quin volum té?

\ex{3} Una piscina fa 25 m de llarg, 10 m d'ample i 2 m de fondària. a) Quants metres quadrats de rajoles calen per a cobrir el fons i les parets? b) Quants litres d'aigua hi caben? c) Quant de temps tardarà a omplir-se amb una bomba que aboca 400 L per minut?

\ex{3} Tres pilotes de tennis de 3,25 cm de radi es venen en un tub cilíndric on caben justes (radi 3,25 cm i altura 19,5 cm). Quina fracció del tub ocupen les pilotes? Quin volum queda buit?

\ex{3} Un got cilíndric de 4 cm de radi té aigua fins a 10 cm d'altura. Hi tirem 5 boles d'1 cm de radi, que queden cobertes. Fins a quina altura arriba ara l'aigua?

\vfill
\begin{center}\color{muted}\small\faCheckCircle[regular]\ \ Ànim! Si et costa algun exercici, repassa els apunts interactius de la unitat.\end{center}

\end{document}
